
\documentclass{article}
\usepackage{colortbl}
\usepackage{makecell}
\usepackage{multirow}
\usepackage{supertabular}

\begin{document}

\newcounter{utterance}

\centering \large Interaction Transcript for game `wordle', experiment `high\_frequency\_words\_no\_clue\_no\_critic', episode 5 with claude{-}opus{-}4{-}8{-}azure.
\vspace{24pt}

{ \footnotesize  \setcounter{utterance}{1}
\setlength{\tabcolsep}{0pt}
\begin{supertabular}{c@{$\;$}|p{.15\linewidth}@{}p{.15\linewidth}p{.15\linewidth}p{.15\linewidth}p{.15\linewidth}p{.15\linewidth}}
    \# & \multicolumn{2}{c}{Player} && \multicolumn{2}{c}{Game Master} \\
    \hline

    \theutterance \stepcounter{utterance}  
    & & & \multicolumn{4}{p{0.6\linewidth}}{
        \cellcolor[rgb]{0.9,0.9,0.9}{
            \makecell[{{p{\linewidth}}}]{
                \texttt{\tiny{[P1$\langle$GM]}}
                \texttt{Ste jezikovni čarovnik, ki rad ugiba besede z uporabo danih pravil.} \\
\\ 
\texttt{Dobrodošli v Wordlu! Imate šest poskusov, da uganete ciljno besedo, veljavno slovensko besedo s petimi malimi črkami. Prosimo, uporabite oznaki "razlaga:" in "ugibanje:", da za vsak poskus podate jedrnato razlago.} \\
\\ 
\texttt{Na primer, če je vaš poskus "otrok", naj bo vaš odgovor} \\
\texttt{razlaga: to je pogosta petčrkovna slovenska beseda in svoj poskus začenjam s to besedo.} \\
\texttt{ugibanje: otrok} \\
\\ 
\texttt{Po vsakem poskusu bo vaš odgovor preverjen in prejeli boste povratne informacije, ki bodo pokazale, katere črke so pravilne (zelena), katere črke so pravilne, vendar na napačnem mestu (rumena), in katere črke so napačne (rdeča). Te povratne informacije so lahko koristne pri določanju, katere črke vključiti ali izključiti v naslednjem poskusu.} \\
\\ 
\texttt{Na primer, povratne informacije za "otrok" bi lahko bile:} \\
\texttt{povratna\_informacija\_ugibanja: o<rumena> t<rumena> r<zelena> o<rumena> k<rdeča>} \\
\\ 
\texttt{Razlaga naj vsebuje podrobnosti o tem, kako se povratna\_informacija\_ugibanja uporabi za oblikovanje novega poskusa.} \\
\\ 
\texttt{Ne ustvarjajte nobenega drugega besedila.} \\
\\ 
\texttt{Začnimo z vašim prvim poskusom.} \\
            }
        }
    }
     \\ \\

    \theutterance \stepcounter{utterance}  
    & \multicolumn{4}{p{0.6\linewidth}}{
        \cellcolor[rgb]{0.561,0.737,0.561}{
            \makecell[{{p{\linewidth}}}]{
                \texttt{\tiny{[P1$\rangle$GM]}}
                \texttt{razlaga: to je pogosta petčrkovna slovenska beseda in svoj poskus začenjam z njo, saj vsebuje več samoglasnikov.} \\
\texttt{ugibanje: mesto} \\
            }
        }
    }
    & & \\ \\

    \theutterance \stepcounter{utterance}  
    & & & \multicolumn{2}{p{0.3\linewidth}}{
        \cellcolor[rgb]{0.9,0.9,0.9}{
            \makecell[{{p{\linewidth}}}]{
                \texttt{\tiny{[GM$|$GM]}}
                \texttt{attempts = 1} \\
\texttt{target\_word = pojem} \\
\texttt{guess = mesto} \\
\texttt{guess\_feedback = m<yellow> e<yellow> s<red> t<red> o<yellow>} \\
            }
        }
    }
    & & \\ \\

    \theutterance \stepcounter{utterance}  
    & & & \multicolumn{4}{p{0.6\linewidth}}{
        \cellcolor[rgb]{0.9,0.9,0.9}{
            \makecell[{{p{\linewidth}}}]{
                \texttt{\tiny{[P1$\langle$GM]}}
                \texttt{povratna\_informacija\_ugibanja: m<rumena> e<rumena> s<rdeča> t<rdeča> o<rumena>} \\
\\ 
\texttt{Odgovorite samo v tej obliki.} \\
\texttt{razlaga: podrobnosti} \\
\texttt{ugibanje: beseda} \\
            }
        }
    }
     \\ \\

    \theutterance \stepcounter{utterance}  
    & \multicolumn{4}{p{0.6\linewidth}}{
        \cellcolor[rgb]{0.561,0.737,0.561}{
            \makecell[{{p{\linewidth}}}]{
                \texttt{\tiny{[P1$\rangle$GM]}}
                \texttt{razlaga: m, e in o so v besedi, a na napačnih mestih; s in t sta izločena. Poskusim "omela", ki vsebuje vse tri črke na drugih mestih.} \\
\texttt{ugibanje: omela} \\
            }
        }
    }
    & & \\ \\

    \theutterance \stepcounter{utterance}  
    & & & \multicolumn{2}{p{0.3\linewidth}}{
        \cellcolor[rgb]{0.9,0.9,0.9}{
            \makecell[{{p{\linewidth}}}]{
                \texttt{\tiny{[GM$|$GM]}}
                \texttt{Error: The guessed word is not a valid word for this game.} \\
            }
        }
    }
    & & \\ \\

    \theutterance \stepcounter{utterance}  
    & & & \multicolumn{4}{p{0.6\linewidth}}{
        \cellcolor[rgb]{0.9,0.9,0.9}{
            \makecell[{{p{\linewidth}}}]{
                \texttt{\tiny{[P1$\langle$GM]}}
                \texttt{Vaše ugibanje ni veljavna beseda za to igro. Poskusite znova.} \\
\\ 
\texttt{Odgovorite samo v tej obliki.} \\
\texttt{razlaga: podrobnosti} \\
\texttt{ugibanje: beseda} \\
            }
        }
    }
     \\ \\

    \theutterance \stepcounter{utterance}  
    & \multicolumn{4}{p{0.6\linewidth}}{
        \cellcolor[rgb]{0.561,0.737,0.561}{
            \makecell[{{p{\linewidth}}}]{
                \texttt{\tiny{[P1$\rangle$GM]}}
                \texttt{razlaga: potrebujem m, e in o na novih mestih (m ne na 1, e ne na 2, o ne na 5). Beseda "pomen" ima o na 2., m na 3. in e na 4. mestu.} \\
\texttt{ugibanje: pomen} \\
            }
        }
    }
    & & \\ \\

    \theutterance \stepcounter{utterance}  
    & & & \multicolumn{2}{p{0.3\linewidth}}{
        \cellcolor[rgb]{0.9,0.9,0.9}{
            \makecell[{{p{\linewidth}}}]{
                \texttt{\tiny{[GM$|$GM]}}
                \texttt{attempts = 2} \\
\texttt{target\_word = pojem} \\
\texttt{guess = pomen} \\
\texttt{guess\_feedback = p<green> o<green> m<yellow> e<green> n<red>} \\
            }
        }
    }
    & & \\ \\

    \theutterance \stepcounter{utterance}  
    & & & \multicolumn{4}{p{0.6\linewidth}}{
        \cellcolor[rgb]{0.9,0.9,0.9}{
            \makecell[{{p{\linewidth}}}]{
                \texttt{\tiny{[P1$\langle$GM]}}
                \texttt{povratna\_informacija\_ugibanja: p<zelena> o<zelena> m<rumena> e<zelena> n<rdeča>} \\
\\ 
\texttt{Odgovorite samo v tej obliki.} \\
\texttt{razlaga: podrobnosti} \\
\texttt{ugibanje: beseda} \\
            }
        }
    }
     \\ \\

    \theutterance \stepcounter{utterance}  
    & \multicolumn{4}{p{0.6\linewidth}}{
        \cellcolor[rgb]{0.561,0.737,0.561}{
            \makecell[{{p{\linewidth}}}]{
                \texttt{\tiny{[P1$\rangle$GM]}}
                \texttt{razlaga: p, o, e so zeleni na mestih 1, 2, 4; m je v besedi, a ne na 3. mestu, torej gre na 5. mesto. "pojem" ustreza vzorcu p{-}o{-}\_{-}e{-}m.} \\
\texttt{ugibanje: pojem} \\
            }
        }
    }
    & & \\ \\

    \theutterance \stepcounter{utterance}  
    & & & \multicolumn{2}{p{0.3\linewidth}}{
        \cellcolor[rgb]{0.9,0.9,0.9}{
            \makecell[{{p{\linewidth}}}]{
                \texttt{\tiny{[GM$|$GM]}}
                \texttt{attempts = 3} \\
\texttt{target\_word = pojem} \\
\texttt{guess = pojem} \\
\texttt{guess\_feedback = p<green> o<green> j<green> e<green> m<green>} \\
            }
        }
    }
    & & \\ \\

    \theutterance \stepcounter{utterance}  
    & & & \multicolumn{2}{p{0.3\linewidth}}{
        \cellcolor[rgb]{0.9,0.9,0.9}{
            \makecell[{{p{\linewidth}}}]{
                \texttt{\tiny{[GM$|$GM]}}
                \texttt{game\_result = WIN} \\
            }
        }
    }
    & & \\ \\

\end{supertabular}
}

\end{document}
